% GENERATED FILE. Edit canonical lessons, then run npm run generate:books.
% canonical-source-hash: fnv1a64:c85fcef2a0c77285
% canonical-lessons: GU-C334-khajano, GU-C334-ajgar, GU-C334-lobh, GU-C334-lalcavvu, GU-C334-bhatakvu

\chapter{The Hidden Treasure}
\label{ch:gu-the-hidden-treasure}
\hlchaptermodality{334}

\begin{chapteropening}
\textbf{By the end of this chapter:} \emph{I can tell what greed did to a man who wanted treasure.}

Complete the last lesson of chapter 334: I can tell what greed did to a man who wanted treasure.
\end{chapteropening}

\section[\texorpdfstring{\gu{ખજાનો} (khajāno)}{ખજાનો (khajāno)}]{\gu{ખજાનો} (khajāno) — a treasure}

\label{lesson:GU-C334-khajano}



\begin{quote}
Before the new one: say the Gujarati for vanished, then the Gujarati for to appear.
\end{quote}

\subsection*{You'll want to know: \gu{ખજાનો}}
\textbf{\gu{ખજાનો}} — \emph{khajāno} — \textquotedblleft{}a treasure\textquotedblright{}.

\begin{quote}
\gu{ગુફામાં} \gu{એક} \gu{જૂનો} \gu{ખજાનો} \gu{હતો}.
\end{quote}

\begin{quote}
\emph{guphāmā̃ ek jūno khajāno hato.}
\end{quote}

\begin{quote}
There was an old treasure in the cave.
\end{quote}

A masculine noun.

\begin{cousinweb}
\textbf{\gu{ખજાનો}} (\emph{khajāno}) is usually traced, through Persian, to Arabic \emph{khizāna}, \textquotedblleft{}a storehouse, a treasury\textquotedblright{}. It is masculine: \textbf{\gu{જૂનો} \gu{ખજાનો}}.
\end{cousinweb}

\begin{grammarlens}[title={What you've built}]
One more word for telling what happened.
\end{grammarlens}

\subsection*{Your turn}
\begin{itemize}
\raggedright
  \item \emph{Say it:} \emph{khajāno}
  \item \emph{Say it:} \emph{khajāno}, once more
  \item \emph{Say it:} \emph{guphāmā̃ ek jūno khajāno hato.} — There was an old treasure in the cave.
  \item \emph{Recall:} say the Gujarati for vanished, then the Gujarati for to appear, then say \emph{khajāno} again
\end{itemize}

\begin{tcolorbox}[breakable,colback=teal!4,colframe=teal!35!black,title={Before you move on}]
What does \gu{ખજાનો} mean? (\textquotedblleft{}A treasure\textquotedblright{}.) Say it once more.
\end{tcolorbox}
\section[\texorpdfstring{\gu{અજગર} (ajgar)}{અજગર (ajgar)}]{\gu{અજગર} (ajgar) — a python}

\label{lesson:GU-C334-ajgar}



\begin{quote}
Before the new one: say the Gujarati for to appear, then the Gujarati for a treasure.
\end{quote}

\subsection*{You'll want to know: \gu{અજગર}}
\textbf{\gu{અજગર}} — \emph{ajgar} — \textquotedblleft{}a python\textquotedblright{}.

\begin{quote}
\gu{ખજાના} \gu{પાસે} \gu{એક} \gu{મોટો} \gu{અજગર} \gu{બેઠો} \gu{હતો}.
\end{quote}

\begin{quote}
\emph{khajānā pāse ek moṭo ajgar beṭho hato.}
\end{quote}

\begin{quote}
A huge python lay beside the treasure.
\end{quote}

A masculine noun.

\begin{cousinweb}
\textbf{\gu{અજગર}} (\emph{ajgar}) is a tatsama, Sanskrit \emph{ajagara}, \textquotedblleft{}a goat-swallower\textquotedblright{}, a large snake (Monier-Williams). In stories it guards treasure.
\end{cousinweb}

\begin{grammarlens}[title={What you've built}]
One more word for telling what happened.
\end{grammarlens}

\subsection*{Your turn}
\begin{itemize}
\raggedright
  \item \emph{Say it:} \emph{ajgar}
  \item \emph{Say it:} \emph{ajgar}, once more
  \item \emph{Say it:} \emph{khajānā pāse ek moṭo ajgar beṭho hato.} — A huge python lay beside the treasure.
  \item \emph{Recall:} say the Gujarati for to appear, then the Gujarati for a treasure, then say \emph{ajgar} again
\end{itemize}

\begin{tcolorbox}[breakable,colback=teal!4,colframe=teal!35!black,title={Before you move on}]
What does \gu{અજગર} mean? (\textquotedblleft{}A python\textquotedblright{}.) Say it once more.
\end{tcolorbox}
\section[\texorpdfstring{\gu{લોભ} (lobh)}{લોભ (lobh)}]{\gu{લોભ} (lobh) — greed}

\label{lesson:GU-C334-lobh}



\begin{quote}
Before the new one: say the Gujarati for a treasure, then the Gujarati for a python.
\end{quote}

\subsection*{You'll want to know: \gu{લોભ}}
\textbf{\gu{લોભ}} — \emph{lobh} — \textquotedblleft{}greed\textquotedblright{}.

\begin{quote}
\gu{લોભને} \gu{લીધે} \gu{ચોર} \gu{ફસાઈ} \gu{ગયો}.
\end{quote}

\begin{quote}
\emph{lobhne līdhe cor phasāī gayo.}
\end{quote}

\begin{quote}
Greed got the thief into trouble.
\end{quote}

A masculine noun.

\begin{cousinweb}
\textbf{\gu{લોભ}} (\emph{lobh}) is a tatsama, Sanskrit \emph{lobha}, \textquotedblleft{}eager desire, greed\textquotedblright{} (Monier-Williams): wanting more than one needs, the fault of many story villains.
\end{cousinweb}

\begin{grammarlens}[title={What you've built}]
One more word for telling what happened.
\end{grammarlens}

\subsection*{Your turn}
\begin{itemize}
\raggedright
  \item \emph{Say it:} \emph{lobh}
  \item \emph{Say it:} \emph{lobh}, once more
  \item \emph{Say it:} \emph{lobhne līdhe cor phasāī gayo.} — Greed got the thief into trouble.
  \item \emph{Recall:} say the Gujarati for a treasure, then the Gujarati for a python, then say \emph{lobh} again
\end{itemize}

\begin{tcolorbox}[breakable,colback=teal!4,colframe=teal!35!black,title={Before you move on}]
What does \gu{લોભ} mean? (\textquotedblleft{}Greed\textquotedblright{}.) Say it once more.
\end{tcolorbox}
\section[\texorpdfstring{\gu{લલચાવવું} (lalcāvvũ)}{લલચાવવું (lalcāvvũ)}]{\gu{લલચાવવું} (lalcāvvũ) — to tempt}

\label{lesson:GU-C334-lalcavvu}



\begin{quote}
Before the new one: say the Gujarati for a python, then the Gujarati for greed.
\end{quote}

\subsection*{You'll want to know: \gu{લલચાવવું}}
\textbf{\gu{લલચાવવું}} — \emph{lalcāvvũ} — \textquotedblleft{}to tempt\textquotedblright{}.

\begin{quote}
\gu{ખજાનાએ} \gu{ચોરને} \gu{લલચાવ્યો}.
\end{quote}

\begin{quote}
\emph{khajānāe corne lalcāvyo.}
\end{quote}

\begin{quote}
The treasure tempted the thief.
\end{quote}

Past, with \gu{તેણે} (he or she …): \gu{તેણે} \gu{બાળકને} \gu{લલચાવ્યું}, \emph{teṇe bāḷakne lalcāvyũ}.

\begin{cousinweb}
\textbf{\gu{લલચાવવું}} (\emph{lalcāvvũ}) is probably built on \emph{lālac}, a longing: CDIAL sets the Hindi cousins \emph{lālac}, \textquotedblleft{}cupidity\textquotedblright{}, and \emph{lalcānā}, \textquotedblleft{}to lust after\textquotedblright{}, under Sanskrit \emph{lālitya}, \textquotedblleft{}amorousness\textquotedblright{} (CDIAL 11029). To tempt is to make someone long for something. The one tempted takes \textbf{\gu{ને}}.
\end{cousinweb}

\begin{grammarlens}[title={What you've built}]
One more word for telling what happened.
\end{grammarlens}

\subsection*{Your turn}
\begin{itemize}
\raggedright
  \item \emph{Say it:} \emph{lalcāvvũ}
  \item \emph{Say it:} \emph{lalcāvvũ}, once more
  \item \emph{Say it:} \emph{khajānāe corne lalcāvyo.} — The treasure tempted the thief.
  \item \emph{Recall:} say the Gujarati for a python, then the Gujarati for greed, then say \emph{lalcāvvũ} again
\end{itemize}

\begin{tcolorbox}[breakable,colback=teal!4,colframe=teal!35!black,title={Before you move on}]
What does \gu{લલચાવવું} mean? (\textquotedblleft{}To tempt\textquotedblright{}.) Say it once more.
\end{tcolorbox}
\section[\texorpdfstring{\gu{ભટકવું} (bhaṭakvũ)}{ભટકવું (bhaṭakvũ)}]{\gu{ભટકવું} (bhaṭakvũ) — to wander (lost)}

\label{lesson:GU-C334-bhatakvu}



\begin{quote}
Before the new one: say the Gujarati for greed, then the Gujarati for to tempt.
\end{quote}

\subsection*{You'll want to know: \gu{ભટકવું}}
\textbf{\gu{ભટકવું}} — \emph{bhaṭakvũ} — \textquotedblleft{}to wander (lost)\textquotedblright{}.

\begin{quote}
\gu{ચોર} \gu{ઘણા} \gu{દિવસ} \gu{જંગલમાં} \gu{ભટક્યો}.
\end{quote}

\begin{quote}
\emph{cor ghaṇā divas jangalmā̃ bhaṭakyo.}
\end{quote}

\begin{quote}
The thief wandered in the forest for many days.
\end{quote}

Past (he …): \gu{તે} \gu{જંગલમાં} \gu{ભટક્યો}, \emph{te jangalmā̃ bhaṭakyo}.

\begin{cousinweb}
is listed by CDIAL under \emph{*bhaṭ}, \textquotedblleft{}sudden movement or noise\textquotedblright{}, glossed \textquotedblleft{}to blaze, start\textquotedblright{} (CDIAL 9365), while CDIAL files the Hindi cousin \emph{bhaṭaknā}, \textquotedblleft{}to wander\textquotedblright{}, under Sanskrit \emph{bhraṣṭá}, \textquotedblleft{}fallen\textquotedblright{} (CDIAL 9655); which line the wandering sense belongs to is not settled. \textbf{\gu{રખડવું}} is to roam; \textbf{\gu{ભટકવું}} is to wander without finding the way.
\end{cousinweb}

\begin{grammarlens}[title={What you've built}]
One more word for telling what happened.
\end{grammarlens}

\subsection*{Your turn}
\begin{itemize}
\raggedright
  \item \emph{Say it:} \emph{bhaṭakvũ}
  \item \emph{Say it:} \emph{bhaṭakvũ}, once more
  \item \emph{Say it:} \emph{cor ghaṇā divas jangalmā̃ bhaṭakyo.} — The thief wandered in the forest for many days.
  \item \emph{Recall:} say the Gujarati for greed, then the Gujarati for to tempt, then say \emph{bhaṭakvũ} again
\end{itemize}

\begin{tcolorbox}[breakable,colback=teal!4,colframe=teal!35!black,title={Before you move on}]
What does \gu{ભટકવું} mean? (\textquotedblleft{}To wander (lost)\textquotedblright{}.) Say it once more.
\end{tcolorbox}
